\documentclass[11pt,a4paper]{article}
\usepackage[T1]{fontenc}
\usepackage[utf8]{inputenc}
\usepackage[english]{babel}
\usepackage{amsmath,amssymb}
\usepackage[margin=2cm]{geometry}
\setlength{\parindent}{0pt}
\setlength{\emergencystretch}{3em}
\allowdisplaybreaks

\title{About \texttt{false\_function\_metrics.tex}}
\author{}
\date{}
\begin{document}
\maketitle
This note explains what the test data in \texttt{false\_function\_metrics.tex} is used for. That file is input for the engine in \texttt{main.c}; it is read by the test \texttt{test\_index\_identities.c} in this folder.

\begin{enumerate}
\item \textbf{Purpose.} The file holds 24 false statements with a diagonal metric whose entries are functions of the coordinates. It is a regression test against false proofs: the engine must declare each one false, with a counterexample certified by the independent checker.

\item \textbf{How the statements were built.} Each one is a true statement of \texttt{function\_metrics.tex} with one small change: an exponent of $r$ or a lost factor; a raised index read as a lowered one, or the metric taken for its inverse or for the identity; a sign. Their falsity was confirmed by the oracle of \texttt{tools/fuzz/fuzz\_tensors.py} in its \texttt{data} mode.

\item \textbf{What the engine needs in order to pass.} The same reading as for the true statements, and a point where the two sides differ. The two values are written only when they are polynomials. The test also alters each certificate and requires the checker to reject it.

\item \textbf{Format expected by the test.} One \texttt{\textbackslash item} per statement, written as \texttt{\$\textbackslash displaystyle ...\$}. The metric is declared as $g_{ij} = \operatorname{diag}(1, r^{2})$: an entry is a number, a polynomial in letters without indices, or a quotient $\frac{\dots}{\dots}$ of two of them. Only there a number above a letter is an exponent; in the tables, in the lists and in the identity a power is written as a product, $r r$.

\item \textbf{How to run it.} From the top folder of the repository:
\begin{verbatim}
cc -std=c11 -O2 -Wall -Wextra -Werror -o /tmp/test_index_identities \
    tests/tensor_calculus/test_index_identities.c && /tmp/test_index_identities
\end{verbatim}

\item \textbf{Expected result.} The line \texttt{PASS: 24 of 24 statements refuted, 24 with a certificate the checker accepts}, and a line \texttt{PASS} for the altered certificates.

\end{enumerate}
\end{document}
